\documentclass[english,french,table,aspectratio=43]{beamer}

\pdfinclusioncopyfonts=1
\pdfgentounicode=1
\pdfminorversion=5 
\pdfcompresslevel=9
\pdfobjcompresslevel=2

%%%%%%%%%%%%%%%%%%%%%%%%%%%%%%%%%%%%%%%%%%%%%%%%%%%%%%%%%%%%%%%%%%%%%
% Packages :
\usepackage{etex}
\usepackage[utf8]{inputenc}          % Attention encodage => UTF 8
\usepackage[T1]{fontenc}             % Codage de caractères T1
\usepackage{lmodern}
\usepackage{cmap}

\rmfamily

% Declare that Latin Modern Roman (lmr) should take
% its bold (b) and bold extended (bx) weight, and small capital (sc) shape, 
% from the corresponding Computer Modern Roman (cmr) font, for the T1 font encoding.
\DeclareFontShape{T1}{lmr}{b}{sc}{<->ssub*cmr/bx/sc}{}
\DeclareFontShape{T1}{lmr}{bx}{sc}{<->ssub*cmr/bx/sc}{}

\usepackage{xspace}
\usepackage{babel}
\usepackage[fixlanguage]{babelbib}
\usepackage[autolanguage]{numprint}
\usepackage[sectionbib]{bibunits}
\usepackage{hyphenat}
\usepackage{csquotes}
\usepackage{subfiles}
\usepackage{amsmath}
\usepackage{amssymb}
\usepackage{amsthm}
\usepackage{mathtools}
\usepackage{latexsym}

\usepackage{circledsteps}
\usepackage{orcidlink}
\usepackage{adjustbox}
%\usepackage{ragged2e}
\usepackage{graphicx}
\usepackage{hyperref}
\usepackage{multirow}
\usepackage{booktabs}
\usepackage{bookmark}
\usepackage{listings}
\usepackage{siunitx}
\usepackage{framed}
\usepackage{cancel}
\usepackage{color}
\usepackage{url}
\usepackage{doi}

\usepackage{tikz}
\usepackage{tikz-dependency}
%\usepackage{tikz-qtree}
%\usepackage{tikz-qtree-compat}
\usepackage{forest}

\usepackage{pgfplots}
\usepackage{pgfplotstable}

\usepackage[svgnames]{xcolor}
%\usepackage[inline]{enumitem}
\usepackage[acronym,shortcuts,translate=babel]{glossaries}
\usepackage[linesnumbered,ruled,vlined,boxed,frenchkw,onelanguage]{algorithm2e}
\usepackage[type={CC},modifier={by-nc-nd},lang=french,version={4.0}]{doclicense}

%\usepackage[a-3u,ua,x-4]{pdfx}

% Bib
\selectbiblanguage{french}

% Algorithm2e
\renewcommand{\CommentSty}[1]{\footnotesize\ttfamily\textcolor{Green}{#1}}
\renewcommand{\KwSty}[1]{\textnormal{\textcolor{blue}{\bfseries #1}}\unskip}
\renewcommand{\FuncSty}[1]{\textsf{#1}}
\renewcommand{\ProcFnt}[1]{\textsf{#1}}

\NewCommandCopy{\legacyunderscore}{\_}
\renewcommand{\_}{\ifincsname_\else\legacyunderscore\fi}

\makeatletter
\def\@algocf@pre@ruled{\begin{shaded}}%
\def\@algocf@post@ruled{\end{shaded}\relax}%
\makeatother

\colorlet{shadecolor}{WhiteSmoke}

% Tikz
\usetikzlibrary{
  automata, graphs, positioning, fit, tikzmark,
  shapes, shapes.multipart, arrows.meta, quotes,
  overlay-beamer-styles
}
\usetikzmarklibrary{listings}

% \drawBrace[xshift]{beginningNode}{endingNode}{comment}
% This command draws a brace between two tikzmarks, to their right, no matter which
% one is the rightmost, and includes a node midway the brace, to write the comment.
\newcommand*{\drawBrace}[4][.5em]{%
    \node[draw=none, fit={(#2) (#3)}] (rectg) {};
    \draw [decoration={brace,amplitude=0.5em},decorate,thick] ([xshift=#1]rectg.north east) -- ([xshift=#1]rectg.south east) node[midway,right,xshift=1em]{#4};
}%

\newcommand*{\drawCircle}[1]{%
    \raisebox{-.5ex}{\tikz{\draw[fill=#1,line width=.5pt] circle[radius=1ex];}}%
}%

\tikzset{
  labeled node/.style={
    rectangle split, rectangle split parts=2,
    rectangle split draw splits=false,
    rounded corners,
    very thick,
    draw=blue!50,
    rectangle split part fill={blue!50, white},
    align=left,
    font=\sffamily\bfseries\boldmath\small
  },
  labeled edge/.style={
    rectangle split, rectangle split parts=1,
    rectangle split draw splits=false,
    rounded corners,
    very thick,
    draw=orange!60,
    rectangle split part fill={orange!60, white},
    align=left,
    font=\sffamily\bfseries\boldmath\small
  },
  labeled property edge/.style={
    rectangle split, rectangle split parts=2,
    rectangle split draw splits=false,
    rounded corners,
    very thick,
    draw=orange!60,
    rectangle split part fill={orange!60, white},
    align=left,
    font=\sffamily\bfseries\boldmath\small
  },
  every two node part/.style={font=\sffamily\small},
}

\tikzstyle{flow-node} = [text centered, minimum height=2em, minimum width=4em, text width=4em, draw]
\tikzstyle{flow-start} = [rectangle, rounded corners, flow-node]
\tikzstyle{flow-process} = [rectangle, text centered, minimum height=2em, text width=7em, draw]
\tikzstyle{flow-data} = [trapezium, trapezium left angle=70, trapezium right angle=110, flow-node]
\tikzstyle{flow-decision} = [diamond, aspect=2.5, flow-node]
\tikzstyle{flow-arrow} = [thick, ->, >=latex, rounded corners]

% Pgf
\usepgfplotslibrary{units}
\usepgfplotslibrary{colorbrewer}
\usepgfplotslibrary{groupplots}

\pgfplotsset{
    compat=newest,
    cycle list/Paired,
    scaled x ticks=false,
    scaled y ticks=false
}

% Listing
\renewcommand{\lstlistingname}{Extrait de code}
\renewcommand{\lstlistlistingname}{Liste des extraits de code}

\lstset{
  commentstyle=\color{Green},
  keywordstyle=\color{blue}\bfseries,
  stringstyle=\color{teal},
  backgroundcolor=\color{WhiteSmoke},
  numberstyle=\footnotesize\color{gray},
  basicstyle=\ttfamily\small,
  breakatwhitespace=true,
  breaklines=true,
  captionpos=b,
  keepspaces=true,
	stepnumber=1,
  numbers=left,
  numbersep=5pt,
  showspaces=false,
  showstringspaces=false,
  showtabs=false,
  tabsize=2,
  frame=lines
}

\lstdefinelanguage{cypher}{
	morekeywords={
		MATCH, OPTIONAL, WHERE, NOT, AND, OR, XOR, RETURN, DISTINCT, ORDER, BY,
    ASC, ASCENDING, DESC, DESCENDING, UNWIND, AS, UNION, WITH, ALL, CREATE,
    DELETE, DETACH, REMOVE, SET, MERGE, SET, SKIP, LIMIT, IN, CALL, CASE, WHEN, EXISTS,
    INDEX, DROP, UNIQUE, CONSTRAINT, EXPLAIN, PROFILE, START, FOREACH, GROUP, HAVING},
	sensitive=true,
	morecomment=[l]{//},
	morecomment=[s]{/*}{*/},
	morestring=[b]{"},
	literate=*
	  {<<}{\color{instruction}\guillemotleft{}}{1}
	  {>>}{\textcolor{instruction}{\guillemotright{}}\color{black}}{1}
}

\lstdefinelanguage{sparql}{
	morekeywords={SELECT, DISTINCT, WHERE, OPTIONAL, FILTER, NOT, EXISTS, MINUS, sameTerm, bound},
}

\mode<presentation> {
    \usetheme{Madrid}
    \useoutertheme[subsection=false]{miniframes}
    \setbeamertemplate{frametitle}[default][center]
    \setbeamertemplate{navigation symbols}{}
}

\makeglossaries          
\loadglsentries{../these/glossaries}

\date{1er Juillet, 2024}
\title[Construction automatique de BD m\'edicale]{Construction automatique de bases de données pour le domaine m\'edical: Int\'egration de texte et maintien de la coh\'erence}
\author[Nicolas Hiot]{
    Nicolas Hiot\inst{1,2}
}
\institute[]{
    \inst{1}LIFO -- Universit\'e d'Orl\'eans, INSA CVL, France\qquad
    \inst{2}EnnovLabs -- Ennov, France\\
    \url{nicolas.hiot@etu.univ-orleans.fr}\qquad
    \url{nhiot@ennov.com}
}
\titlegraphic{
    \includegraphics[height=1cm]{../logos/logoOrleans2022.png}
    \hspace{.8cm}
    \includegraphics[height=1cm]{../logos/logoLIFO.png}
    \hspace{.8cm}
    \includegraphics[height=.9cm]{../logos/logoEnnov.png}
    \hspace{.8cm}
    \includegraphics[height=.8cm]{../logos/logoDOING.png}
}

% \AtBeginSection{
%   \begin{frame}{Sommaire}
%     \tableofcontents[currentsection]
%   \end{frame}
% }

\begin{document}

\frame{\titlepage}

\section{Introduction}

% \begin{frame}{Contexte}
%   Textes médicaux : cas cliniques, pharmacovigilance, cosméto-vigilance, etc.\pause
%   \begin{description}
%       \item[Objectif] Exploiter ces données en les rendant interrogeables\pause
%       \item[Solution] Construire automatiquement une base de données\pause
%           \begin{itemize}
%               \item Structurer automatiquement les information du texte
%               \item Garantir la qualité et la cohérence des données
%           \end{itemize}
%   \end{description}
% \end{frame}

\begin{frame}{Données textuelles médicales}
  \begin{itemize}
      \item Le domaine médical génère en quantité des données textuelles difficiles à exploiter efficacement.
      \begin{itemize}
          \item Cas cliniques
          \item Pharmacovigilance
          \item Cosmetovigilance
      \end{itemize}
      \vfill
      \item Décrivent l'historique des patients:
      \begin{itemize}
          \item Diagnostiques
          \item Traitements et réactions
          \item Antécédents médicaux
      \end{itemize}
      \vfill
      \item Importance pour la prise de décision clinique:
      \begin{itemize}
          \item Meilleure compréhension des cas complexes
          \item Amélioration de la qualité des soins
          \item Détection précoce des effets indésirables
      \end{itemize}
  \end{itemize}
\end{frame}

\begin{frame}{Axes de Recherches}
  \begin{itemize}
      \item Proposition : Construire automatiquement une base de données à partir d'un corpus textuel médical
      \vfill\pause
      \item Objectifs :
      \begin{itemize}
          \item Faciliter l'accès à des informations précises et structurées
          \item Permettre l'interrogation et l'exploitation efficaces des données
      \end{itemize}
      \vfill\pause
      \item Axes de Recherches :
      \begin{itemize}
          \item Extraction et structuration automatique de l'information
          \item Assurer la qualité et la cohérence des données:
      \end{itemize}
  \end{itemize}
\end{frame}

% \begin{frame}{Contexte}
%   \begin{block}{Genèse}
%     Textes médicaux : cas cliniques, pharmacovigilance, cosméto-vigilance, etc.\pause
%     \begin{description}
%         \item[Objectif] Exploiter ces données en les rendant interrogeables\pause
%         \item[Solution] Construire une base de données
%             \begin{itemize}
%                 \item Structurer les information pour les rendre exploitable
%                 \item Garantir la qualité et la cohérence des données
%             \end{itemize}
%     \end{description}
%     \end{block}

%   \pause

%   \begin{block}{Contraintes}
%     \begin{itemize}
%         \item Données spécifiques à chaque client
%         \item Absence de données annotées
%         \item Puissance de calcul et quantité de données limitées
%         \item Nécessité d'un système rapide et efficace
%         %\item Besoin d'explicabilité des résultats
%     \end{itemize}
%     \end{block}
% \end{frame}

% \begin{frame}{Corpus médical}
%     \begin{block}{Corpus de cas cliniques en français (DEFT)} % \cite{grabar2020cas} 
%         \begin{itemize}
%             \item 100 fichiers dé-identifiés
%             \item Situations réelles et fictives
%             \item Imbrication d'entités
%             \item Annoté  en entités (8350 annotations), mais pas en relations
%         \end{itemize}
%     \end{block}

%     \begin{block}{Entités}
%         \begin{itemize}
%             \item Patient (anatomie)
%             \item Pratique clinique (examen, pathologie, signe ou symptôme)
%             \item Traitements (dose, durée, fréquence, mode, substance, traitement, valeur)
%             \item Temporalité (date, moment)
%         \end{itemize}
%     \end{block}
% \end{frame}

% \begin{frame}{Relations}
%     \begin{table}[htb]
%         \centering
%         \footnotesize
%         \resizebox{\textwidth}{!}{ 
%             \def\arraystretch{1.1}%
%             \begin{tabular}{p{2.2cm}p{3.3cm}p{8cm}}
%                 %\hline
%                 Relation                             & Paire d'entités                               & Exemples                                                                                                                                                                 \\ \hline\hline
%                 \multirow{2}{*}{\emph{means}}        & \multirow{2}{*}{\emph{mode}-\emph{substance}} & \multirow{2}{*}{\shortstack[l]{\emph{une} [\emph{crème}]{\textsuperscript{\emph{mode}}} \emph{de} [\emph{nistatin}]{\textsuperscript{\emph{sub}}} \emph{a été préscrite} \\
%                 'an [ointment]{\textsuperscript{\emph{mode}}} of [nistatin]{\textsuperscript{\emph{sub}}} was prescribed'}}                                                                                                                                                     \\
%                                                      &                                               &                                                                                                                                                                          \\
%                 \hline

%                 \multirow{2}{*}{\emph{measure\_of}}  & \emph{dose}-\emph{substance}                  & \multirow{2}{*}{\shortstack[l]{\emph{la} [\emph{doxorubicine}]{\textsuperscript{\emph{sub}}} \emph{à raison de} [\emph{37,5 mg/m2/dose}]{\textsuperscript{\emph{dose}}}  \\ '[\emph{37,5 mg/m2/dose}]{\textsuperscript{\emph{dose}}} of [doxorubicine]{\textsuperscript{\emph{sub}}}'}}  \\
%                                                      & \emph{value}-\emph{substance}                 &                                                                                                                                                                          \\
%                 \hline

%                 \multirow{4}{*}{\emph{accompagnies}} & \emph{pathology}-\emph{pathology}             & \multirow{4}{*}{\shortstack[l]{\emph{le patient présentait une} [\emph{fatigue importante}]{\textsuperscript{\emph{sosy}}} \emph{de même}                                \\\emph{qu'une} [\emph{hyperthermie}]{\textsuperscript{\emph{sosy}}} \\ 'the patient displayed [great tiredness]{\textsuperscript{\emph{sosy}}} as well \\ as [hyperthermia]{\textsuperscript{\emph{sosy}}}'}} \\
%                                                      & \emph{sosy}-\emph{sosy}                       &                                                                                                                                                                          \\
%                                                      & \emph{pathology}-\emph{sosy}                  &                                                                                                                                                                          \\
%                                                      & \emph{substance}-\emph{substance}             &                                                                                                                                                                          \\
%                 \hline

%                 \multirow{3}{*}{\shortstack[l]{\emph{reveals}/                                                                                                                                                                                                                  \\\emph{searches}/\\\emph{tests}}}  & \emph{exam}-\emph{value} & \multirow{3}{*}{\shortstack[l]{\emph{l'}[\emph{échographie abdominale et pelvienne}]{\textsuperscript{\emph{exam}}} \emph{révèle la} \\ \emph{présence d'une} [\emph{masse surrénalienne à droite}]{\textsuperscript{\emph{sosy}}} \\ 'the [abdominal and pelvic ultrasonography]{\textsuperscript{\emph{exam}}} reveals \\ a [growth in the right adrenal gland]{\textsuperscript{\emph{sosy}}}' }} \\
%                                                      & \emph{exam}-\emph{sosy}                       &                                                                                                                                                                          \\
%                                                      & \emph{exam}-\emph{pathology}                  &                                                                                                                                                                          \\
%                                                      &                                               &                                                                                                                                                                          \\
%                 \hline

%                 \multirow{4}{*}{\emph{location}}     & \emph{anatomy}-\emph{pathology}               & \multirow{4}{*}{\shortstack[l]{\emph{elle a subi une} [\emph{résection au niveau}                                                                                        \\ \emph{du} [\emph{lobe supérieur droit}]{\textsuperscript{\emph{anat}}}]{\textsuperscript{\emph{treat}}} \\ 'she underwent a [resection of the [upper right lobe]{\textsuperscript{\emph{anat}}}]{\textsuperscript{\emph{treat}}}'}} \\
%                                                      & \emph{anatomy}-\emph{sosy}                    &                                                                                                                                                                          \\
%                                                      & \emph{anatomy}-\emph{exam}                    &                                                                                                                                                                          \\
%                                                      & \emph{anatomy}-\emph{treatment}               &                                                                                                                                                                          \\
%                 %\hline
%             \end{tabular}
%         }
%     \end{table}
% \end{frame}

%\part{Intégration de données textuelles}

\subfile{parts/struct}

%\part{Maintenance et cohérence d'une base de données incomplète}

\subfile{parts/nulls}

% \part{Extraction d'informations}

% \frame{\partpage}

%\subfile{parts/tal}

\section{Conlusion}

\begin{frame}{}
    \setbeamercolor{block body}{bg=blue, fg=white}
    \begin{block}{}
        \Huge
        \centering
        \vspace{1em}
        Conclusion et perspectives
        \vspace{1em}
    \end{block}
\end{frame}

\begin{frame}{Contributions}
    \begin{block}{Structuration automatique de l'information}
        \begin{itemize}
            \item Méthode combinant approches symboliques et statistiques pour structurer les données
            \item Formalisation du processus de structuration avec des règles de réécriture d'arbres
            \item Proposition d'une métagrammaire correspondant à un modèle général (indépendante du modèle de données)
        \end{itemize}
    \end{block}
    \begin{block}{Qualité et cohérence des données}
        \begin{itemize}
            \item Politique de mise à jour incrémentale pour le maintient de la cohérence d'une base de données
            \item Approche basée sur un \acrshort{sgbd} permettant le passage à l'echelle
            \item Utilisation de bases de données graphes et de requêtes récursives pour récupérer les LinkedNulls
        \end{itemize}
    \end{block}
\end{frame}

\begin{frame}{Productions scientifiques et perspectives}
    \begin{columns}
        \begin{column}{.5\linewidth}
            7 publications scientifiques :
            \begin{itemize}
                \item 1 revue (IJIWIS)
                \item 2 conférences internationales (iiWAS, ADBIS)
                \item 4 workshops (DEFT, Doing)
            \end{itemize}
        \end{column}
        \begin{column}{.5\linewidth}
            Participation projets :
            \begin{itemize}
                \item ANR SENDUP
                \item APR-IA Doing
                \item MaDICS - Action Doing
            \end{itemize}
        \end{column}
    \end{columns}
    2 logiciels :
    \begin{itemize}
        \begin{columns}
            \begin{column}{.5\linewidth}
                \item ArchiTxt : Structuration automatique
            \end{column}
            \begin{column}{.5\linewidth}
                \item DataFix : Maintenance cohérente
            \end{column}
        \end{columns}
    \end{itemize}
    \vfill
    \begin{block}{Perspectives}
        \begin{itemize} 
            \item Intégrer les contraintes dès la structuration automatique (guider la structuration pour garantir leur respect)
            \vfill
            \item Créer une vue unifiée de multiples sources de données (IA, data journalisme, etc.)
        \end{itemize}
    \end{block}
\end{frame}

\frame{\titlepage}

% \begin{frame}[allowframebreaks]{References}
%     \bibliographystyle{apalike}
%     \bibliography{../these/biblio}
% \end{frame}

\end{document}
